\scanpage{30}
\margin{\handwritten{SGA 4 $\overline{\text{\uline{XVI}}}$}}
\mspage{VII 2}

Du point de vue substance, 2.1.1 se ramène trivialement au cas où dans $X\handwritten{\longrightarrow}S$, $X$ est affine, donc quasi-compact quasi-séparé (on dit maintenant ``cohérent''), et alors c'est une conséquence du th. de changement de base par $f$, qui est donné dans SGA 4 XVI 1.1\inserted{et qui est essentiellement}. Mais dans loc. cit. on utilise toute la force du théorème de changement de base propre, ce qui est en effet abusif, car en interprétant comme un énoncé d'acyclicité globale 2.1.1, on est ramené comme tu l'observes \deleted{à 2.1.5} au cas particulier 2.1.5, qui se ramène trivialement à son tour au cas du degré de transcendance 1, et par là au th. de changement de\inserted{corps de} base pour le $H^1$ premier à $p$ pour une courbe propre et lisse. Je pense présenter les choses ainsi en réédition de SGA 4, de sorte que ton 2.1.1 y sera énoncé en toutes lettres, s'il ne l'est déjà (je n'ai pas les textes en main). Cela te dispensera donc de le prouver, ta seule tâche consistant à interpréter en termes de champs cet énoncé \deleted{(que tu pourrais donner sous forme de rappels)}.
\margin{\figure{\begin{tikzpicture}\draw (0,0) .. controls (.45,-.12) and (.60,-.15) .. (1.1,.5);\end{tikzpicture}}}

Du point de vue rédaction, on a l'impression que tu t'es ingénié à tout présenter à l'\unsure{e|E}nvers. Il faut en tous cas bloquer ensemble 2.1.1 et 2.1.7 (pratiquement identiques) et 2.1.5 (celui-ci en corollaire, peut-être) \deleted{- ce sera dans les rappels}\inserted{\handwritten{cf texte joint}}\deleted{. Avant même de donner ces énoncés, il faudra avoir rappelé les définitions, pour que ça ait un sens.} \deleted{Au lieu de pro-lisse resp. lisse, il suffit\inserted{dans le tH.} de supposer tout \uline{localement} \uline{$L$-\inserted{\unsure{1-|l-}}asphérique} (``par exemple ...'' ).} D'autre part, \deleted{les énoncés} 2.12, 2.1.3, 2.1.4, 2.1.7.1 n'ont rien à voir avec les schémas, il sont manifestement importants en eux-mêmes et devraient passer dans un numéro à part dans \handwritten{\uline{IV}\deleted{\unsure{6|5}}\inserted{ou \uline{V}}} (avec un peu d'ordre dans la suite des énoncés)
\margin{\handwritten{\unsure{si on utilise} \unsure{ailleurs} la ``factorisation canonique d'un champ'',}}

\deleted{Enfin, les compléments schématiques 2.1.7.2 à 2.1.7.4 \unsure{de|d}}
% archive sheet number: 29
